% =====================================================================
\section{Primitives, Indexed Rendering and Memory}\label{sec:prims}
\pipestrip{1}
Every primitive is a list of unique vertices $(x,y,z,n_x,n_y,n_z)$, 24 bytes each, plus a triangle index list, wound counter-clockwise seen from outside so that \code{GL\_CCW} and back-face culling keep only outward faces. The resolution is $N=36$ slices, $\Delta\phi=10^\circ$. The generators in \code{src/primitives/} are the \emph{only} source of geometry, except the terrain disc (the cylinder cap subdivided into rings and lifted by a height field).

\subsection{Cube, sphere, cylinder, cone and frustum}
\textbf{Cube.} Each face has an outward normal $\vect n$ and in-plane axes $\vect u,\vect v$ with
$\vect u\times\vect v=\vect n$; its corners are
$\vect p_c=\tfrac12\vect n+\tfrac12 s_u\vect u+\tfrac12 s_v\vect v$ with
$(s_u,s_v)\in\{(-,-),(+,-),(+,+),(-,+)\}$, counter-clockwise about $\vect n$, split into $(0,1,2)$ and
$(0,2,3)$. A corner is stored once per face (its normal differs): 24 vertices, 36 indices, 12 triangles (T01--T03).

\textbf{Sphere.} With $r=0.5$, $\theta_k=k\pi/S$, $\phi_j=2\pi j/N$ ($S=18$, $N=36$):
$x=r\sin\theta_k\cos\phi_j$, $y=r\cos\theta_k$, $z=r\sin\theta_k\sin\phi_j$, $\vect n=\vect p/r$. Only the
northern hemisphere is computed; the southern half is its reflection $y\mapsto-y$. A reflection reverses orientation, so each mirrored triangle
$(a,b,c)$ is emitted as $(a',c',b')$. The shared equator gives $2+(S-1)N=614$ vertices and $2N(S-1)=1224$
triangles (T04--T06).

\textbf{Cylinder, cone, frustum.} The axis is the local $z$ axis, $z\in[-0.5,0.5]$. The side joins two rings,
\begin{equation}
  \vect p^{\,b}_j=(r_b\cos\phi_j,\,r_b\sin\phi_j,\,-0.5),\quad
  \vect p^{\,t}_j=(r_t\cos\phi_j,\,r_t\sin\phi_j,\,+0.5),\quad
  \vect n_j=\mathrm{normalize}(\cos\phi_j,\,\sin\phi_j,\,r_b-r_t),
\end{equation}
and each cap is a fan around its own centre. With $r_b=r_t=0.5$ this is the lab cylinder: $4N+2=146$ vertices,
$4N=144$ triangles. With $r_t=0$ the same routine yields the \textbf{cone} (109 vertices, 72 triangles; no degenerate triangles, T36); with $0<r_t<r_b$ and no caps it yields the
searchlight \textbf{frustum} (72 vertices, T10). Cones (pine whorls, missile noses; 259 instances in the export) thus need no extra generator.

\begin{figure}[h]
\centering
\picsheet{sheet_vertex_cube}{\linewidth}
\caption{\sheetcaption{sheet_vertex_cube}}
\end{figure}
\begin{figure}[h]
\centering
\picdia{A1_cube_vertices}{0.92\linewidth}
\caption{The cube: 24 vertices (4 per face with the face normal) and 36 indices; 720\,B indexed against 864\,B flat.}
\end{figure}
\begin{figure}[h]
\centering
\picdia{A2_sphere_rings}{0.92\linewidth}
\caption{The sphere: rings of 36 slices, normals equal to the normalised positions, the southern half mirrored.}
\end{figure}
\begin{figure}[h]
\centering
\picdia{A3_cylinder_cone}{0.92\linewidth}
\caption{Cylinder, cone and frustum from one generator.}
\end{figure}
\begin{figure}[h]
\centering
\picdia{A4_terrain_disc}{0.92\linewidth}
\caption{The terrain disc: ring structure (left) and the real triangles with the entrance-ramp hole (right).}
\end{figure}
\subsection{The generator code}
Listing~\ref{lst:Cube}: each face is a frame $(\vect n,\vect u,\vect v)$ with four counter-clockwise corners. Listing~\ref{lst:Sphere}: the northern hemisphere is generated ring by ring and mirrored. Listing~\ref{lst:Cylinder}: two rings joined by side quads with normals tilted by $r_b-r_t$, plus caps; a top radius of 0 gives the cone.
\lstCube
\lstSphere
\lstCylinder
\paragraph{The terrain disc.} The terrain is the cylinder cap subdivided into $R=130$ rings of $N=384$ slices ($1+RN=49\,921$ vertices, $N(2R-1)=99\,456$ triangles minus 2\,016 cut out for the entrance trench $=97\,440$, T13), radius 36\,m, each vertex lifted to
$y=h(x,z)$: rolling hills, a levelled bunker pad, a creek, a track, a berm and a mound over the bunker. Normals come from central differences of $h$.

\subsection{Indexed rendering and memory}
\sheetfig{s02_cube_indices}
\sheetfig{s06_geometry_stats}
\sheetfig{s06_charts}
Without indexing every triangle stores its three vertices (lab technique 1: 18 floats for a square); with
indexing, unique vertices go to the VBO and triangles reference them through the EBO (technique 2: 12 floats $+$ 6
indices). The VAO records \emph{location 0: offset 0, stride 24\,B} (position) and \emph{location 1: offset 12\,B}
(normal); one VAO+VBO+EBO per primitive is shared by every part. Memory per mesh is
$M_\text{flat}=3T\cdot 24$\,B and $M_\text{indexed}=24V+12T$\,B (Table~\ref{tab:memory}). The six meshes take \memIndexed\,B indexed against \memNonIndexed\,B non-indexed, dominated by the
terrain disc; they serve all \nLeaves{} parts. T12 checks that every indexed mesh is smaller than its flat version.

\begin{table}[h]
\centering\small
\caption{GPU meshes (from \code{build/scene\_export.json}). The terrain's vertex count (49\,921) is that of the 36\,m mesh (T13).}
\label{tab:memory}
% generated by report/gen_tables.py from build/scene_export.json -- do not edit
\begin{tabular}{@{}lrrrrrrr@{}}
\toprule
Mesh & $V$ & Indices & Triangles & Non-indexed (B) & Indexed (B) & Saving & Instances\\\midrule
Cube & 24 & 36 & 12 & 864 & 720 & 16.7\,\% & 2,231\\
Sphere & 614 & 3,672 & 1,224 & 88,128 & 29,424 & 66.6\,\% & 816\\
Cylinder & 146 & 432 & 144 & 10,368 & 5,232 & 49.5\,\% & 2,068\\
Cylinder (frustum) & 72 & 216 & 72 & 5,184 & 2,592 & 50.0\,\% & 3\\
Terrain disc & 49,921 & 292,320 & 97,440 & 7,015,680 & 2,367,384 & 66.3\,\% & 1\\
Cone & 109 & 216 & 72 & 5,184 & 3,480 & 32.9\,\% & 295\\
\midrule
Sum & & & & 7,125,408 & 2,408,832 & 66.2\,\% &\\
\bottomrule
\end{tabular}

\end{table}

\begin{figure}[h]
\centering
\picsheet{sheet_memory}{\linewidth}
\caption{\sheetcaption{sheet_memory}}
\end{figure}
\paragraph{Instancing and culling.} The grass is the unit cube drawn with \code{glDrawElementsInstanced} (per-instance buffer, \code{glVertexAttribDivisor}); the quality preset sets the count: LOW 14\,000, HIGH 35\,000,
ULTRA 140\,000, EPIC 450\,000 blades. Rain streaks, foliage cards and effect sprites use the same mechanism. Every pass culls parts whose bounding sphere $r=\tfrac12\sqrt{s_x^2+s_y^2+s_z^2}$ lies outside one of the six
frustum planes (Gribb--Hartmann, Chapter~\ref{sec:geomchain}).

\paragraph{Tests.} T01--T13 check counts, centring, normal length, winding, sphere reflection, cylinder
axis, index validity and the memory comparison (all \nTests{} tests: Section~\ref{sec:testing}).
